\section*{Capítulo \ref{ch:g11:lithosphere-asthenosphere} --- Litosfera y astenosfera}

\begin{solution}{exo:g11:lithosphere-asthenosphere:1}
La \omterm{def:g11:lithosphere-asthenosphere:lithosphere}{litosfera} es la capa externa rígida de la Tierra, formada por la corteza y
la parte más alta del manto; la \omterm{def:g11:lithosphere-asthenosphere:lithosphere}{astenosfera} es el manto más caliente y más
débil que hay debajo, que fluye lentamente a lo largo de mucho tiempo. La
corteza oceánica, la corteza continental y la parte más alta del manto forman
parte de la \omterm{def:g11:lithosphere-asthenosphere:lithosphere}{litosfera}; el núcleo externo, no.
\end{solution}

\begin{solution}{exo:g11:lithosphere-asthenosphere:2}
El iceberg se hunde hasta apartar su propio peso de agua: la fracción
sumergida es $910/1025 \approx 0.888$, alrededor del \qty{89}{\percent}. La
fracción que asoma es de alrededor del \qty{11.2}{\percent}, así que la altura
total es de $20/0.112 \approx \qty{180}{m}$, de los que unos \qty{160}{m}
están bajo el agua.
\end{solution}

\begin{solution}{exo:g11:lithosphere-asthenosphere:3}
$15 + 26 \times 3 = \qty{93}{\celsius}$. A \qty{100}{km}, el mismo gradiente
daría más de \qty{2600}{\celsius}, muy por encima de la temperatura de fusión
de la roca del manto; pero las \omterm{def:g11:seismic-structure:body-waves}{ondas S} atraviesan el manto superior, que por
tanto es sólido. El gradiente tiene que hacerse mucho menor con la
profundidad.
\end{solution}

\begin{solution}{exo:g11:lithosphere-asthenosphere:4}
La franja de 0 a \qty{500}{m} sobre el nivel del mar, con alrededor del
\qty{16}{\percent} de la superficie, y la franja de 4000 a \qty{4500}{m} bajo
el nivel del mar, con alrededor del \qty{11}{\percent} (la franja de 4500 a
\qty{5000}{m} por debajo es casi igual de grande).
\end{solution}

\begin{solution}{exo:g11:lithosphere-asthenosphere:5}
Con $P = 0$: \qty{1121}{\celsius}. A \qty{3}{GPa}: $-5.104 \times 9 +
132.899 \times 3 + 1120.661 \approx -45.9 + 398.7 + 1120.7 \approx
\qty{1473}{\celsius}$.
\end{solution}

\begin{solution}{exo:g11:lithosphere-asthenosphere:6}
Unos \qty{730}{\celsius} a \qty{50}{km} (\omterm{def:g11:lithosphere-asthenosphere:lithosphere}{litosfera}), unos
\qty{1370}{\celsius} a \qty{150}{km} y \qty{1410}{\celsius} a
\qty{250}{km} (\omterm{def:g11:lithosphere-asthenosphere:lithosphere}{astenosfera}). La temperatura sube más entre 50 y
\qty{150}{km}, más de \qty{600}{\celsius}, porque la mayor parte de ese
intervalo está en la \omterm{def:g11:lithosphere-asthenosphere:lithosphere}{litosfera}; entre 150 y \qty{250}{km} solo sube unos
\qty{45}{\celsius}.
\end{solution}

\begin{solution}{exo:g11:lithosphere-asthenosphere:7}
Raíz: $3 \times 2830/550 \approx \qty{15.4}{km}$. Espesor adicional:
$3 \times 3380/550 \approx \qty{18.4}{km}$, es decir, los \qty{3}{km} de la
cordillera más su raíz.
\end{solution}

\begin{solution}{exo:g11:lithosphere-asthenosphere:8}
La corteza es más gruesa en $55 - 41.1 = \qty{13.9}{km}$, y $\Delta H = h
\times 3380/550$, luego $h = 13.9 \times 550/3380 \approx \qty{2.3}{km}$.
\end{solution}

\begin{solution}{exo:g11:lithosphere-asthenosphere:9}
Disminución: $(4.470 - 4.419)/4.470 \approx \qty{1.1}{\percent}$. Salto:
$(4.644 - 4.419)/4.419 \approx \qty{5.1}{\percent}$. Con la profundidad, la
roca está cada vez más comprimida, lo que hace más rápidas las ondas
sísmicas; una disminución significa que otra cosa se impone: la roca está
caliente, cerca de la fusión, y es más débil.
\end{solution}

\begin{solution}{exo:g11:lithosphere-asthenosphere:10}
(1) $(1345 - 610)/60 \approx \qty{12}{\celsius}$ por kilómetro. (2) Es más o
menos la mitad del gradiente en la superficie. Parte del calor se produce
dentro de la corteza superior; ese calor también tiene que atravesar las
capas altas para escapar, lo que hace más fuerte el gradiente cerca de la
superficie. Bajo la corteza, solo atraviesa la roca el calor que llega de más
abajo, y el gradiente es menor.
\end{solution}

\begin{solution}{exo:g11:lithosphere-asthenosphere:11}
A \qty{200}{km}: $1300 + 90 = \qty{1390}{\celsius}$; a \qty{400}{km}:
$1300 + 180 = \qty{1480}{\celsius}$. La subida es de \qty{90}{\celsius} en
\qty{200}{km}. A través de la \omterm{def:g11:lithosphere-asthenosphere:lithosphere}{litosfera}, la temperatura sube
$1345 - 15 = \qty{1330}{\celsius}$ en solo \qty{100}{km}: unas treinta veces
más por kilómetro.
\end{solution}

\begin{solution}{exo:g11:lithosphere-asthenosphere:12}
(1) $118 + 12.3\,z = 1300 + 0.45\,z$ da $z = 1182/11.85 \approx
\qty{100}{km}$. (2) $118 + 21.5\,z = 1300 + 0.45\,z$ da $z = 1182/21.05
\approx \qty{56}{km}$. (3) El oeste, más caliente, tiene una \omterm{def:g11:lithosphere-asthenosphere:lithosphere}{litosfera} mucho
más delgada, como se mide (unos \qty{40}{km} donde es más delgada); la
\omterm{def:g11:lithosphere-asthenosphere:lithosphere}{litosfera} muy gruesa del centro del continente necesita un gradiente todavía
menor.
\end{solution}

\begin{solution}{exo:g11:lithosphere-asthenosphere:13}
(1) Por la proposición, $\Delta H = h\,\rho_m/(\rho_m - \rho_c)$, luego $h =
\Delta H\,(\rho_m - \rho_c)/\rho_m$. (2) $550/3380 \approx 0.16$: alrededor de
una sexta parte de cualquier corteza adicional aparece como altura, y cinco
sextas partes como raíz. (3) En el iceberg, la fracción que sobresale del
agua es $(1025 - 910)/1025 \approx 0.11$: la misma regla, (densidad del
fluido $-$ densidad del cuerpo que flota)/densidad del fluido.
\end{solution}

\begin{solution}{exo:g11:lithosphere-asthenosphere:14}
(1) $1489 - 1345 = \qty{144}{\celsius}$. (2) Un poco de fundido se reúne en
películas finas entre los granos de la roca. Los granos dejan entonces de
estar firmemente trabados: la roca resiste peor la cizalla, así que las \omterm{def:g11:seismic-structure:body-waves}{ondas
S} se frenan, y puede deformarse y fluir con mucha más facilidad.
\end{solution}

\begin{solution}{exo:g11:lithosphere-asthenosphere:15}
(1) No. Aplanar las tierras emergidas adelgazaría un poco la corteza
continental, pero seguiría siendo gruesa y ligera, y seguiría flotando mucho
más alta que la corteza delgada y densa del fondo oceánico; los océanos
seguirían llenando las cuencas entre los continentes. (2) Hay dos tipos de
corteza, de espesor y densidad distintos: una corteza continental gruesa y
ligera y una corteza oceánica delgada y densa.
\end{solution}

\begin{solution}{pb:g11:lithosphere-asthenosphere:1}
\textbf{1.} $5.01 - 0.67 = \qty{4.34}{km}$.

\textbf{2.} $(3380 - 2830)/3380 \approx \qty{16}{\percent}$.

\textbf{3.} Quedaría por debajo la fracción $2830/3380 \approx 0.84$; el
iceberg se hunde algo más, $910/1025 \approx 0.89$, porque la densidad del
hielo está más cerca de la del agua que la de la corteza de la del manto.

\textbf{4.} Durante los segundos de una onda sísmica, o el tiempo de un
terremoto, el manto es un sólido; solo a lo largo de muchísimo tiempo fluye
la \omterm{def:g11:lithosphere-asthenosphere:lithosphere}{astenosfera} y deja flotar la corteza.

\textbf{5.} Por debajo del \omterm{def:g11:seismic-structure:moho}{Moho} más profundo, las dos columnas están hechas
de la misma roca del manto a las mismas profundidades, y allí pesan lo mismo;
todas las diferencias están por encima de esa profundidad y hay que
incluirlas.

\textbf{6.} El continente medio: corteza desde \qty{0.67}{km} por encima
hasta \qty{40.4}{km} por debajo del nivel del mar, y luego manto hasta
\qty{70}{km}. La meseta: corteza desde \qty{5.01}{km} sobre el nivel del mar
hasta un \omterm{def:g11:seismic-structure:moho}{Moho} más profundo en el valor de la raíz, y luego manto hasta
\qty{70}{km}.

\textbf{7.} La columna de la meseta tiene $h + r$ más de corteza; la del
continente tiene, en cambio, manto en el espesor $r$: $\rho_c\,(h + r) = \rho_m\,r$,
luego $r = h\,\rho_c/(\rho_m - \rho_c)$ y $\Delta H = h + r =
h\,\rho_m/(\rho_m - \rho_c)$.

\textbf{8.} $3380/550 \approx 6.15$.

\textbf{9.} $\Delta H = 4.34 \times 6.15 \approx \qty{26.7}{km}$; la raíz es
$26.7 - 4.34 \approx \qty{22.3}{km}$.

\textbf{10.} El \omterm{def:g11:seismic-structure:moho}{Moho} del continente medio está a $41.1 - 0.67 \approx
\qty{40.4}{km}$ bajo el nivel del mar; el de la meseta, a $40.4 + 22.3 \approx
\qty{62.8}{km}$ bajo el nivel del mar.

\textbf{11.} $41.1 + 26.7 \approx \qty{67.8}{km}$.

\textbf{12.} Un kilómetro de roca sobre un metro cuadrado tiene una masa de
$\rho \times \qty{1000}{m^3}$: \num{2830} toneladas de corteza o \num{3380}
toneladas de manto. Continente: $41.1 \times 2830 + (70 - 40.43) \times 3380
\approx \num{116300} + \num{99900} \approx \num{216300}$ toneladas. Meseta:
$67.77 \times 2830 + (70 - 62.76) \times 3380 \approx \num{191800} +
\num{24500} \approx \num{216300}$ toneladas. Las dos son iguales.

\textbf{13.} $41.1 - 0.67 \times 6.15 \approx \qty{37.0}{km}$.

\textbf{14.} $67.8/41.1 \approx 1.65$ veces.

\textbf{15.} $3300/470 \approx 7.02$, $\Delta H \approx 4.34 \times 7.02
\approx \qty{30.5}{km}$, y una corteza de unos \qty{71.6}{km}. Un cambio de
alrededor del \qty{2}{\percent} en la densidad del manto cambia el espesor en
unos \qty{4}{km}: el resultado depende de forma apreciable de las densidades,
pero se mantiene cerca de \qty{70}{km}.

\textbf{16.} El modelo da unos \qty{68}{km}, y las ondas sísmicas unos
\qty{70}{km}: concuerdan con un margen de pocos puntos porcentuales.

\textbf{17.} $h = 3.23 - 0.67 = \qty{2.56}{km}$, $\Delta H = 2.56 \times
6.15 \approx \qty{15.7}{km}$, y una corteza de unos \qty{56.8}{km}, unos
\qty{7}{km} más de lo medido.

\textbf{18.} Más delgada. Una columna cuya parte alta es más ligera que la
media pesa menos con el mismo espesor, así que flota más alta: una corteza
más ligera necesita una raíz menor para estar a una altura dada. Por tanto,
la corteza real puede ser más delgada que la del modelo, como se mide.

\textbf{19.} A lo largo de mucho tiempo, la \omterm{def:g11:lithosphere-asthenosphere:lithosphere}{astenosfera} fluye: se aparta bajo
una columna pesada y vuelve bajo una ligera, hasta que las columnas se
equilibran.

\textbf{20.} \textbf{La corteza bajo el Tíbet tiene unos \qty{70}{km} de
espesor}, mucho más que los \qty{41}{km} de un continente medio: la meseta
está a cinco kilómetros de altura porque su corteza, gruesa y ligera, flota
sobre el manto, más denso, con una raíz de más de \qty{20}{km}.
\end{solution}
